% ============================================================
%  ch2_background_concepts.tex
%  Chapter 2: Background and Concepts
%  Master's Thesis -- SQL MATCH_RECOGNIZE on Pandas DataFrames
% ============================================================
\chapter{Background and Concepts}
\label{chap:background}

This chapter establishes the theoretical and practical foundations
required to understand the design and implementation of the
\texttt{MATCH\_RECOGNIZE} engine for Pandas DataFrames.
Section~\ref{sec:bg-pattern-matching} introduces the problem of
sequential pattern detection in relational data.
Section~\ref{sec:bg-match-recognize} examines the SQL:2016
\texttt{MATCH\_RECOGNIZE} clause in detail.
Chapter~\ref{chap:methodology} describes the implementation and
performance optimisation techniques.  Chapter~\ref{chap:related}
surveys related work and identifies
the research gap addressed by this thesis.

% ----------------------------------------------------------------
\section{Pattern Matching in Sequential Data Analysis}
\label{sec:bg-pattern-matching}
% ----------------------------------------------------------------

Sequential pattern detection represents a fundamental challenge in
data analysis, where the objective is to identify ordered sequences
of events or states within large datasets.  In contrast to
conventional queries that match rows individually against one or
more predicates, pattern matching takes preceding and subsequent
rows into consideration and returns groups of rows when they follow
a specified pattern.  This makes it possible to retrieve specific
shapes---such as V-, W-, or U-shapes---within the data.  Pattern
matching is therefore well suited to event-driven processes such as
security monitoring, stock-price analysis, sensor-data analysis,
fraud detection, and complex event processing~\cite{OracleMR}.

Unlike traditional relational operations that focus on individual
rows or aggregated values, pattern matching examines the temporal
or logical ordering of data points to extract meaningful sequences.
This capability is essential across many domains including financial
market analysis, fraud detection, IoT sensor monitoring, and
behavioural analytics.

Historically, relational database systems excelled at set-based
operations such as selection, projection, joins, and aggregation,
but lacked expressive mechanisms for sequence analysis.  Analysts
requiring pattern detection were forced to export data to specialised
stream-processing systems, implement procedural logic using cursors
and temporary tables, or employ application-layer processing with
significant performance overhead.  These limitations created a gap
between the widespread availability of sequential data and the tools
available to analyse it effectively.  The introduction of row pattern
recognition in SQL:2016 addresses this deficiency by bringing
sequence analysis capabilities directly into the query
language~\cite{ISO2016SQL}.

% ................................
\subsection{Real-World Applications}
\label{subsec:bg-applications}
% ................................

\paragraph{Financial Market Analysis.}
Detecting V-shaped price patterns---price decline followed by
recovery---requires examining ordered sequences of price movements.
For example, identifying stocks whose price troughs are successively
higher (a ``higher-low'' pattern) signals a potential upward trend.
The canonical example used throughout this thesis is the detection
of such V-patterns in customer order data~\cite{OracleMR}.

\paragraph{Fraud Detection.}
Credit card fraud often manifests as specific transaction sequences,
such as multiple small test charges followed by a large purchase, or
geographically impossible transactions occurring within unrealistic
timeframes.  Expressing these rules declaratively in SQL is more
maintainable than equivalent procedural logic~\cite{petkovic2022fraud}.

\paragraph{IoT and Sensor Data.}
Manufacturing systems monitor sensor readings to detect equipment
degradation patterns, such as a gradual temperature rise followed by
a sudden drop, which may indicate imminent failure.  Pattern queries
over time-series sensor streams are a natural fit for
\texttt{MATCH\_RECOGNIZE}~\cite{OracleMR}.

\paragraph{Clickstream and Funnel Analysis.}
E-commerce platforms analyse user navigation sequences to identify
conversion paths or abandonment patterns.  A classic funnel query
detects the sequence
``view~$\rightarrow$ add-to-cart~$\rightarrow$ remove~$\rightarrow$
exit,'' which signals purchase hesitation.  With
\texttt{MATCH\_RECOGNIZE}, the analyst can also quantify \emph{how
many} add-to-cart events occurred before removal, extract the exact
timestamps of each step, and compute the total dwell time in a single
declarative query---tasks that would otherwise require multiple
correlated subqueries or procedural session-stitching logic.
Such funnel-analysis workloads are common examples of row-pattern
recognition in analytical SQL~\cite{hoffa2021snowflake}.

\paragraph{Network Security and Intrusion Detection.}
Cyber-attacks rarely manifest as a single anomalous event; they
unfold as ordered sequences of actions.  A typical multi-stage
intrusion follows the pattern: port scanning~$\rightarrow$ repeated
failed authentication attempts~$\rightarrow$ a single successful
login~$\rightarrow$ privilege escalation.
\texttt{MATCH\_RECOGNIZE} can express this entire kill-chain as one
pattern query over an event log, flagging the offending session and
extracting the first and last timestamps, the targeted port, and the
account compromised---all without exporting the log to a separate
SIEM tool~\cite{petkovic2022fraud}.

\paragraph{Healthcare and Patient Monitoring.}
In clinical data streams, early warning systems are designed to detect
physiological deterioration before a critical event.  A typical
sepsis-onset pattern consists of rising heart rate rows, followed by
a drop in blood pressure rows, followed by a spike in white-blood-cell
count.  This example illustrates how ordered clinical measurements
could be modelled as a row-pattern query over a time-series of vital
signs, while the exact alerting workflow would remain domain- and
system-dependent.

\paragraph{Supply Chain and Logistics.}
Operational disruptions in supply chains often emerge as a cascade:
a single late shipment~$\rightarrow$ repeated partial deliveries~$\rightarrow$
a stockout event.  \texttt{MATCH\_RECOGNIZE} applied to an order
fulfilment log can identify such escalation sequences per supplier or
warehouse and measure the time between first delay and eventual
stockout, supporting earlier investigation rather than purely
retrospective reporting.

\paragraph{Financial Compliance and Trade Surveillance.}
Trade-surveillance workloads often search for potential market
manipulation patterns, for example ``layering,'' where a trader
places a large resting order
(to move the price), then executes a trade on the opposite side, then
cancels the resting order.  This three-step sequence---place, trade,
cancel---with precise timing constraints is naturally expressed as a
\texttt{MATCH\_RECOGNIZE} pattern, turning a complex compliance rule
into an auditable SQL query.


% ----------------------------------------------------------------
\section{The SQL:2016 \texttt{MATCH\_RECOGNIZE} Clause}
\label{sec:bg-match-recognize}
% ----------------------------------------------------------------

The SQL:2016 standard introduced
\texttt{MATCH\_RECOGNIZE}~\cite{ISO2016SQL,iso2016}, a construct
that enabled row pattern matching in SQL.  Row pattern matching is
a technique that makes it possible to find sequences of rows based on a
pattern defined by regular expressions.  Using a syntax similar to
regular expressions, users can describe and search for complex
sequences of events~\cite{winand2017sqlnew}.  This process facilitates the
recognition of patterns in data series that would otherwise require
the use of additional SQL clauses, such as the \texttt{WINDOW}
clause.  In a theoretical sense, \texttt{MATCH\_RECOGNIZE} does
not add new relational expressiveness: many row-pattern queries can
be rewritten using joins, window functions, recursive queries, or
procedural code.  Its practical contribution is instead a concise
declarative abstraction for sequence-oriented logic that would
otherwise be verbose and difficult to maintain.

At the time of writing, \texttt{MATCH\_RECOGNIZE} is still absent from
several widely used data systems, including
PostgreSQL~\cite{howe2020matchrecognize} and
DuckDB~\cite{findling2025bringing}.  This thesis contributes to
the broader goal of making declarative row-pattern matching available
in environments that lack native support, with a particular focus on
Pandas DataFrames and in-process Python analytics.

The SQL standard distinguishes between different feature levels for
row pattern matching, which can be used to classify
implementations~\cite{winand_match_recognize}.
\textbf{Feature~R010} provides the most basic support and serves
as a basis for further features; it includes row pattern
recognition inside the \texttt{FROM} clause and basic aggregate
functions (\texttt{MIN}, \texttt{MAX}, \texttt{SUM},
\texttt{COUNT}, \texttt{AVG}) inside the \texttt{MEASURES}
clause.  \textbf{Feature~R020} extends the areas of application:
\texttt{MATCH\_RECOGNIZE} can be used within the \texttt{WINDOW}
and \texttt{OVER} clauses to describe windows more precisely.
\textbf{Feature~R030} allows extended aggregate support inside the
\texttt{MATCH\_RECOGNIZE}
clause~\cite{ISO2021,michels2018new,melton2017sqlrpr}.

% ................................
\subsection{Clause Structure and Components}
\label{subsec:bg-clause-structure}
% ................................

The clause operates on a rowset (table or derived table) and is
commonly described through eight interconnected sub-clauses, some
mandatory and some optional, as shown in
Listing~\ref{lst:mr-syntax}.  The \texttt{PATTERN} and
\texttt{DEFINE} clauses form the core mandatory part of row pattern
recognition; the remaining clauses control partitioning, ordering,
output, and overlap behaviour.

\begin{lstlisting}[
    language     = SQL,
    caption      = {General syntax of the \texttt{MATCH\_RECOGNIZE} clause},
    label        = {lst:mr-syntax}
]
SELECT <columns>
FROM   <table>
MATCH_RECOGNIZE (
    [PARTITION BY <partition_columns>] -- 1. Data organisation
    [ORDER BY     <ordering_columns>]  -- 2. Row sequencing
    [MEASURES     <expressions>]       -- 3. Output computations
    [<rows_per_match>]                 -- 4. Result granularity
    [AFTER MATCH SKIP <option>]        -- 5. Overlap handling
    PATTERN ( <pattern_regex> )        -- 6. Sequence specification
    [SUBSET <union_variable_defs>]     -- 7. Union variables
    DEFINE <variable_definitions>      -- 8. Variable predicates
) AS <alias>
\end{lstlisting}

\begin{enumerate}

\item \textbf{PARTITION BY.}
  When supplied, divides the input rowset into independent partitions,
  analogous to window-function partitioning.  Pattern matching runs
  independently within each partition, enabling both parallel
  processing and domain-specific grouping (e.g.\ per customer or per
  sensor).  If omitted, the input is treated as a single partition.

\item \textbf{ORDER BY.}
  Establishes the logical ordering of rows within each partition.
  When supplied, the pattern-matching algorithm processes rows
  strictly according to this ordering; for deterministic sequence
  analysis, an explicit ordering is normally required.  Multiple sort
  keys with \texttt{ASC}/\texttt{DESC} modifiers are supported.

\item \textbf{MEASURES.}
  Optionally defines computed expressions that extract values from
  matched row sequences using pattern variables and navigation functions
  (\texttt{FIRST}, \texttt{LAST}, \texttt{PREV}, \texttt{NEXT}).
  Each \texttt{MEASURE} expression becomes a column in the output.

\item \textbf{Output options.}
  \begin{itemize}
    \item \texttt{ONE ROW PER MATCH}: one summary row per complete
      match.
    \item \texttt{ALL ROWS PER MATCH}: one row for every row
      participating in a match.
    \item \texttt{WITH UNMATCHED ROWS}: an
      \texttt{ALL ROWS PER MATCH} modifier that includes
      non-matching rows in the result.
  \end{itemize}

\item \textbf{AFTER MATCH SKIP.}
  Controls how the engine resumes after finding a match:
  \begin{itemize}
    \item \texttt{SKIP TO NEXT ROW} -- resume from the row after
      match start (maximum overlaps).
    \item \texttt{SKIP PAST LAST ROW} -- resume from the row after
      match end (non-overlapping matches).
    \item \texttt{SKIP TO FIRST} \textit{var} -- resume from the
      first row classified as \textit{var}.
    \item \texttt{SKIP TO LAST} \textit{var} -- resume from the last
      row classified as \textit{var}.
  \end{itemize}

\item \textbf{PATTERN.}
  Specifies the sequence using regular-expression-like syntax with
  pattern variables as terminals:
  \begin{itemize}
    \item Concatenation: \texttt{A~B~C} (A followed by B followed by C).
    \item Alternation: \texttt{A~|~B} (either A or B).
    \item Quantifiers: \texttt{A*} (zero or more), \texttt{A+}
      (one or more), \texttt{A?} (zero or one), \texttt{A\{n\}}
      (exactly $n$), \texttt{A\{n,m\}} (between $n$ and $m$).
    \item Grouping: \texttt{(A~B)+} (one or more AB pairs).
    \item Anchors: \texttt{\^{}A~B~C\$} (pattern must span the
      entire partition).
    \item Exclusions: \texttt{\{-~A+~B+~-\}} (rows matching the
      enclosed sub-pattern are excluded from row-level output).
    \item Permutations: \texttt{PERMUTE(A,~B,~C)} (any ordering
      of A, B, and C).
  \end{itemize}

\item \textbf{SUBSET.}
  An optional clause for defining union variables, each of which
  consists of a set of variables already defined in the
  \texttt{PATTERN} clause.  A union variable does not introduce new
  matching behaviour; it provides a shared name for rows matched by
  any of its component variables, and may appear inside
  \texttt{MEASURES}, \texttt{DEFINE}, and skip expressions where
  supported.

\item \textbf{DEFINE.}
  Assigns Boolean predicates to pattern variables.  During matching,
  candidate rows are evaluated against these predicates in the
  current match context to determine whether they may be classified
  as a given variable.  Predicates may reference current-row
  columns, navigation functions, aggregations over matched
  subsequences, and Boolean/comparison
  operators~\cite{winand_match_recognize}.  A pattern variable with
  no \texttt{DEFINE} predicate acts as an unconstrained variable.

\end{enumerate}

A \texttt{MATCH\_RECOGNIZE} clause is embedded in a
\texttt{SELECT} query and can therefore be further processed by
additional specifications.  Possible extensions include applying
groupings or aggregate functions to the outer
\texttt{SELECT}~\cite{winand_match_recognize}.

% ................................
\subsection{Illustrative Example: Stock Price V-Pattern}
\label{subsec:bg-example}
% ................................

Listing~\ref{lst:mr-example} shows a \texttt{MATCH\_RECOGNIZE}
query that detects V-shaped price patterns in customer order data.

\begin{lstlisting}[
    language = SQL,
    caption  = {Detecting V-shaped price patterns with \texttt{MATCH\_RECOGNIZE}},
    label    = {lst:mr-example}
]
SELECT customer_id,
       start_price, bottom_price, final_price,
       start_date,  final_date
FROM   orders
MATCH_RECOGNIZE (
    PARTITION BY customer_id
    ORDER BY     order_date
    MEASURES
        START.price          AS start_price,
        LAST(DOWN.price)     AS bottom_price,
        LAST(UP.price)       AS final_price,
        START.order_date     AS start_date,
        LAST(UP.order_date)  AS final_date
    ONE ROW PER MATCH
    AFTER MATCH SKIP PAST LAST ROW
    PATTERN (START DOWN+ UP+)
    DEFINE
        DOWN AS price < PREV(price),
        UP   AS price > PREV(price)
);
\end{lstlisting}

The engine partitions rows by \texttt{customer\_id}, sorts by
\texttt{order\_date}, and seeks sequences matching
\texttt{START~DOWN+~UP+}: one \texttt{START} row (any row, as it
carries no \texttt{DEFINE} predicate), one or more \texttt{DOWN}
rows (each with a strictly lower price than its predecessor), and
one or more \texttt{UP} rows (each with a strictly higher price
than its predecessor).  For every match found, the \texttt{MEASURES}
sub-clause extracts the starting price and date from the
\texttt{START} row, the lowest price from the last \texttt{DOWN}
row, and the final price from the last \texttt{UP} row.

Chapter~\ref{chap:related} examines existing implementations in
greater depth, comparing their feature coverage, deployment models,
and practical restrictions.
